% Resume focus: Mechanical
% Tailor bullets toward structures, aerodynamics, manufacturing, CAD,
% analysis, test, and hardware build ownership.
\documentclass[10pt]{article}

\usepackage[margin=0.33in]{geometry}
\usepackage{titlesec}
\usepackage{enumitem}
\usepackage{hyperref}
\usepackage[T1]{fontenc}
\usepackage[scaled]{helvet}
\usepackage[none]{hyphenat}
\sloppy

\hypersetup{hidelinks}
\renewcommand{\familydefault}{\sfdefault}
\setlength{\parindent}{0pt}
\setlength{\parskip}{0pt}
\pagestyle{empty}

\setlength{\parskip}{0pt}
\newcommand{\ressection}[1]{
  \vspace{0.12em}
  {\large\bfseries #1}\par
  \vspace{0.05em}\hrule\vspace{0.12em}
}

\setlist[itemize]{
  leftmargin=1.1em,
  itemsep=-1pt,
  topsep=-1pt,
  parsep=0pt,
  partopsep=0pt,
  noitemsep
}

\begin{document}

\begin{center}
{\LARGE {\bfseries Santhosh Sankar}} \\
{\small Aerospace Engineer with Experience in Mechanical Systems and Mechanisms} \\
Ann Arbor, MI \enspace $\vert$ \enspace (734) 846-5255 \enspace $\vert$ \enspace
\href{mailto:ssankar@umich.edu}{ssankar@umich.edu} \enspace $\vert$ \enspace
\href{https://www.linkedin.com/in/santhosh-sankar-3241521b8/}{linkedin.com/in/Tosh-Sankar} \enspace $\vert$ \enspace
\href{https://github.com/ssankar7775}{github.com/ssankar7775}
\end{center}

\ressection{EDUCATION}

{\bfseries University of Michigan}, Ann Arbor, MI \hfill Jan 2026 -- May 2027 (Expected) \\
Master of Engineering in Space Systems \hfill {\bfseries GPA: 4.0/4.0} \\
Climate and Space Sciences and Engineering Department (CLaSP)

\begin{itemize}
  \item Specialization in Controls and Embedded Systems
  \item Relevant Graduate Coursework: Management of Space Systems, Control System Analysis and Design, Astrodynamics
\end{itemize}

{\bfseries University of Michigan}, Ann Arbor, MI \hfill Aug 2021 -- Aug 2024 \\
Bachelor of Science and Engineering in Aerospace Engineering \\
Department of Aerospace Engineering

\ressection{WORK EXPERIENCE}

{\bfseries SpaceX}, Graduate Engineer -- Payload Integration, Launch Engineering \hfill Aug 2025 -- Dec 2025 \\
Vandenberg, CA
\begin{itemize}
\item Engineered and Deployed an additional payload avionics checkout rack, increasing payload processing throughput by \textbf{50\%} to support increased launch cadence.
\item Developed a simulation environment to validate automated payload avionics checkout sequences, enabling \textbf{hardware-independent testing} and engineer training.
\\item Automated payload checkout test plans using \\textbf{Python-based tooling}, reducing required engineering support by \\textbf{1 hour per launch campaign}.
\item Integrated and verified payload avionics hardware during Falcon 9 launch processing, supporting \textbf{vehicle-payload interface validation} and launch readiness testing.
\end{itemize}

\ressection{PROJECT EXPERIENCE}

{\bfseries Michigan Aeronautical Science Association (MASA)}, Aerodynamics and Recovery Lead \hfill Jun 2022 -- Aug 2023 \\
Ann Arbor, MI
\begin{itemize}
  \item Led aerodynamics and recovery development for the \textbf{23-ft liquid bi-propellant Clementine rocket}, coordinating \textbf{40 engineers across 11 subprojects} from preliminary design through launch integration.
  \item Authored black-powder manufacturing SOPs and a pyrotechnic safety training program approved by facility EHS, \textbf{eliminating uncontrolled ignition risk} across all propulsion and recovery test campaigns.
  \item Redesigned the rocket stage separation system from a pyrotechnic piston to pneumatic architecture in <2 weeks, enabling \textbf{successful system-level ground testing} ahead of launch readiness.
  \item Drove fin and fin-can design using MATLAB, FEA, OpenRocket, and CFD, personally executing \textbf{composite layups and structural fabrication} to meet mass and stiffness targets.
\end{itemize}

{\bfseries CubeSat Flight Lab -- University of Michigan}, Systems Engineer \hfill Aug 2023 -- Aug 2024 \\
Ann Arbor, MI
\begin{itemize}
  \item Led a high-altitude balloon (HAB) flight campaign from concept through launch, producing \textbf{CONOPS, system requirements, and full integration/test procedures} for the MC-10 CubeSat program.
  \item Iterated \textbf{6 FlatSat/StratoSat hardware builds}, reducing architecture from \textbf{3U (7 kg) to 1U (1.2 kg)} while maturing multi-subsystem integration and verification testing.
  \item Designed and fabricated an RP2040-based flight termination unit, {\bfseries reducing subsystem mass by 83\%} and validating operation via \textbf{thermal chamber testing at 110,000-ft simulated altitude}.
\end{itemize}

{\bfseries Human Powered Submarine}, Propulsion and Hull Design Team Member \hfill Sep 2023 -- Aug 2024 \\
Ann Arbor, MI
\begin{itemize}
  \item Migrated team CAD workflow from Dassault tools to \textbf{Siemens NX + Teamcenter}, enabling large-assembly collaboration and standardizing \textbf{GD\&T practices} across the design team.
  \item Developed composite hull manufacturing SOPs reducing production time from \textbf{2 weeks to 5 days}, mitigating pre-competition schedule risk.
  \item Designed a drivetrain tensioner system for gear-and-pulley propulsion, \textbf{eliminating chain skipping} and enabling diver-accessible adjustment during underwater operation.

\end{itemize}

{\bfseries MRacing Formula SAE}, Aerodynamics Team Member \hfill Jan 2022 -- May 2022 \\
Ann Arbor, MI
\begin{itemize}
  \item Performed Star-CCM+ CFD analysis of NACA airfoil geometries to quantify \textbf{downforce/drag trade-offs} and inform front and rear wing design decisions.
  \item Fabricated \textbf{12 carbon-fiber composite wing layups} using wet layup and vacuum bagging, establishing repeatable manufacturing quality for competition hardware.

\end{itemize}

\ressection{TECHNICAL SKILLS}

\textbf{Mechanical Design:} Siemens NX, SolidWorks, Teamcenter, CAD modeling, GD\&T \\
\textbf{Programming \& Analysis:} Python, C/C++, CFD, FEA, MATLAB, Simulink, Star-CCM+, OpenRocket \\
\textbf{Manufacturing \& Test:} Composite fabrication, machining, prototyping, integration, hardware build execution, thermal chambers, wind tunnels, instrumentation, verification planning \\
\textbf{Project Management:} Requirements, test planning, technical documentation, team leadership
\end{document}
